\chapter{Esercizi}
\textbf{Es. 1}\\
Calcolare la lunghezza d'onda di un elettrone con energia cinetica $K_\alpha = \qty{6}{\mega\electronvolt}$.\\
\textbf{Soluzione:}\\
La lunghezza d'onda di un elettrone con energia cinetica $K_\alpha = \qty{6}{\mega\electronvolt}$ può essere calcolata utilizzando la relazione di de Broglie:
\begin{equation}
\lambda = \frac{h}{p}
\end{equation}
Otteniamo la quantità di moto $p$ utilizzando l'approssimazione relativistica $K_\alpha \approx E_\alpha$ e dunque:
\begin{equation}
p = \sqrt{2m_e K_\alpha}
\end{equation}
dove $m_e$ è la massa dell'elettrone. Sostituendo $p$ nella formula della lunghezza d'onda:
\begin{equation}
\lambda = \frac{h}{\sqrt{2m_e K_\alpha}}
\end{equation}
Con i seguenti valori:
\begin{itemize}
 \item $h \approx \qty{6.626e-34}{\joule\second}$
 \item $m_e \approx \qty{9.109e-31}{\kilo\gram}$
 \item $K_\alpha$ espresso in joule (convertendo da MeV con $\qty{1}{\mega\electronvolt} = \qty{1.602e-13}{\joule}$)
\end{itemize}
Sostituendo i valori numerici, si ottiene $\lambda \approx \qty{0.05}{\angstrom}$.\\
Utilizzando la formula semiempirica di massa e assumendo che le forze nucleari siano identiche per tutte le coppie di
nucleoni, possiamo scrivere l'energia di legame per nucleone come:\\
\textbf{Es. 2}\\
Fascio di particelle $\alpha$ con una corrente $i_\alpha = \qty{10}{\nano\ampere}$ colpisce un bersaglio di $^{24}\text{Mg}$ ($23.985 \,u$)con $\rho d = \qty{1}{\milli\gram\per\centi\meter^2}$. Creando una reazione isotropa:
\begin{equation}
^{24}_{12}\text{Mg} + ^4_2\alpha \rightarrow ^{27}_{13}\text{Al} + p
\end{equation}
Otteniamo un angolo solido di $\Delta \Omega = \qty{2e-4}{\steradian}$ con un numero di protoni al secondo $\dot{N}_p = \qty{20}{\second^{-1}}$. Calcolare la sezione d'urto differenziale totale.\\
\textbf{Soluzione:}\\
Sappiamo il numero di protoni che colpiscono il bersaglio per unità di tempo nel rivelatore posizionato ad un angolo solido $\Delta \Omega$, dunque otteniamo il numero totale di protoni prodotti per unità di tempo come:
\begin{equation}
  \dot{N}_{p,tot} = \dot{N}_{p}(\Delta \Omega) \dfrac{4 \pi}{\Delta \Omega} \approx \qty{1.26e6}{\second^{-1}}
\end{equation}
Vogliamo trovare il flusso di particelle $\alpha$ che colpiscono il bersaglio:
\begin{equation}
 \phi_\alpha = \dfrac{i_\alpha}{Aq_\alpha}
\end{equation}
Il numero di centri di scattering è dato dai possibili centri di scattering del bersaglio di alluminio.
Possiamo anche scrivere $N_{\text{Mg}}$ in funzione della densità fornita:
\begin{equation}
 \dfrac{m_{\text{Mg}}N_{\text{Mg}}}{dA}d = \rho d
\end{equation}
Dunque abbiamo che:
\begin{equation}
  \sigma = \dfrac{\dot{N}_{p,tot}}{\phi_\alpha N_{\text{Mg}}} = \dfrac{\dot{N}_{p,tot}}{\dfrac{i_\alpha}{\cancel{A}q_\alpha} \dfrac{\rho d \cancel{A}}{m_{\text{Mg}}}} = \dfrac{\dot{N}_{p,tot} m_{\text{Mg}}q_\alpha}{i_\alpha \rho d A} \approx \qty{160}{\milli\barn}
\end{equation}
\textbf{Es. 3}\\
Un elettrone con quantità di moto $p_-$ urta un positrone che si muove con quantità di moto $p_+$ uguale ed opposta. Nell'urto vengono prodotti una coppia di muoni con carica opposta:
\begin{equation}
 e^- + e^+ \rightarrow \mu^- + \mu^+
\end{equation}
Note le masse:
\begin{itemize}
 \item $m_e = \qty{0.511}{\mega\electronvolt\per c^}$
 \item $m_\mu = \qty{105.7}{\mega\electronvolt\per c^2}$
\end{itemize}
Determinare:
\begin{itemize}
  \item Le quantità di moto dei muoni in funzione di quelle degli elettroni
  \item Il valore minimo della quantità di moto degli elettroni affinché la reazione avvenga
\end{itemize}
\textbf{Soluzione:}\\
Utilizzando la conservazione della quantità di moto e dell'energia, possiamo scrivere:
\begin{equation}
 p_- + p_+ = p_{\mu^-} + p_{\mu^+} = 0
\end{equation}
\begin{equation}
 E_- + E_+ = E_{\mu^-} + E_{\mu^+}
\end{equation}
Dove:
\begin{equation}
 E_- = \sqrt{p_-^2 c^2 + m_e^2 c^4}
\end{equation}
\begin{equation}
 E_+ = \sqrt{p_+^2 c^2 + m_e^2 c^4}
\end{equation}
Dunque $E_- = E_+$ e $E_{\mu^-} = E_{\mu^+}$, quindi:
\begin{equation}
 \begin{split}
  2E_{e^-} &= 2E_{\mu^-} \\[8pt]
  E_{e^-} &= E_{\mu^-}
 \end{split}
\end{equation}
Sapendo che:
\begin{equation}
 E_{\mu^-} = \sqrt{p_{\mu^-}^2 c^2 + m_\mu^2 c^4}
\end{equation}
\begin{equation}
 E_{\mu^+} = \sqrt{p_{\mu^+}^2 c^2 + m_\mu^2 c^4}
\end{equation}
Sostituendo le quantità di moto opposte degli elettroni e risolvendo il sistema di equazioni, otteniamo:
\begin{equation}
 p_{\mu^\pm} = \sqrt{\dfrac{E_{e^-}^2}{c^2} - m_\mu^2 c^2} = \sqrt{p_{e^-}^2 + (m_e^2 - m_\mu^2) c^2}
\end{equation}
Otteniamo la quantità di moto minima iniziale ponendo $p_{\mu^\pm} = 0$:
\begin{equation}
 p_{e^-,min} = c \sqrt{m_\mu^2 - m_e^2} \approx \qty{105.5}{\mega\electronvolt\per c}
\end{equation}
\section{Modello a gas di Fermi}
\textbf{Es. 1}\\
Calcolare l'impulso e l'energia di Fermi dei nucleoni nel nucleo $^{16}_8O$ assumendo una distribuzione a simmetria sferica. L'energia di legame totale del nucleo è $E_B = \qty{128}{\mega\electronvolt}$.\\
Calcolare la profondità della buca di potenziale che lega i nucleoni nel nucleo. (Si trascuri la differenza di massa tra protone e neutrone $M_P \approx M_n = \qty{939}{\mega\electronvolt}$)\\
\textbf{Soluzione:}\\
Abbiamo che:
\begin{equation}
 p^n_F = p^p_F = p_F = \hbar \brackets{3 \pi^2 \dfrac{N}{V}}^{1/3}
\end{equation}
Dove $N = 8$ e $V = \dfrac{4}{3} \pi R^3$ con $R = r_0 A^{1/3}$ e $r_0 \approx \qty{1.2}{\femto\meter}$. Dunque:
\begin{equation}
 p_F = \hbar \brackets{3 \pi^2 \dfrac{N}{\dfrac{4}{3} \pi r_0^3 A}}^{1/3} \approx \qty{250}{\mega\electronvolt\per c}
\end{equation}
Ora troviamo l'energia di Fermi:
\begin{equation}
 \varepsilon_F = \sqrt{p_F^2 c^2 + M_N^2 c^4} - M_N c^2 \approx \qty{33}{\mega\electronvolt}
\end{equation}
Per calcolare la profondità della buca di potenziale sappiamo che:
\begin{equation}
  u = \varepsilon_F + |E_B|/A \approx \qty{41}{\mega\electronvolt}
\end{equation}
\textbf{Es. 2}\\
Supponiamo di osservare sperimentalmente il decadimento di una particella in due fotoni. L'angolo di apertura è $\theta = 32^\circ$ e l'energia di ciascun fotone è $E_{\gamma1} = \qty{200}{\mega\electronvolt}$, $E_{\gamma2} = \qty{300}{\mega\electronvolt}$. Calcolare la massa della particella decaduta e la sua quantità di moto.\\
\textbf{Soluzione:}\\
Utilizzando la conservazione dell'energia e della quantità di moto, possiamo scrivere:
\begin{equation}
 E = E_{\gamma1} + E_{\gamma2} = \qty{500}{\mega\electronvolt}
\end{equation}
\begin{equation}
 \vec{p} = \vec{p}_{\gamma1} + \vec{p}_{\gamma2}
\end{equation}
La quantità di moto dei fotoni è data da:
\begin{equation}
 p_{\gamma1} = \dfrac{E_{\gamma1}}{c}
\end{equation}
\begin{equation}
 p_{\gamma2} = \dfrac{E_{\gamma2}}{c}
\end{equation}
L'energia iniziale della particella decaduta è data dalla relazione di dispersione relativistica:
\begin{equation}
 E^2 = p^2 c^2 + m^2 c^4
\end{equation}
Scegliamo l'asse $x$ lungo la direzione della particella decaduta, dunque:
\begin{equation}
 p_x = p_{\gamma1} \cos(\theta_1) + p_{\gamma2} \cos(\theta_2) = P
\end{equation}
\begin{equation}
 p_y = p_{\gamma1} \sin(\theta_1) - p_{\gamma2} \sin(\theta_2) = 0
\end{equation}
Otteniamo:
\begin{equation}
  \dfrac{\sin(\theta_1)}{\sin(\theta_2)} = \dfrac{p_{\gamma2}}{p_{\gamma1}} = \dfrac{E_{\gamma2}}{E_{\gamma1}} = 1.5
\end{equation}
Sapendo che $\theta = \theta_1 + \theta_2 = 32^\circ$, tramite le formule di addizione e sottrazione degli angoli otteniamo:
\begin{equation}
 \sin(\theta_1) = \sin(32^\circ - \theta_2) = \sin(32^\circ) \cos(\theta_2) - \cos(32^\circ) \sin(\theta_2)
\end{equation}
Da cui:
\begin{equation}
  \tan(\theta_2) = \dfrac{\sin(32^\circ)}{1.5 + \cos(32^\circ)} \approx 0.183
\end{equation}
Dunque $\theta_2 \approx 12.7^\circ$ e $\theta_1 \approx 19.3^\circ$. Ora possiamo calcolare la quantità di moto della particella decaduta:
\begin{equation}
 P = \dfrac{E_{\gamma1}}{c} \cos(19.3^\circ) + \dfrac{E_{\gamma2}}{c} \cos(12.7^\circ) \approx \qty{481}{\mega\electronvolt}
\end{equation}
Ora possiamo calcolare la massa della particella decaduta:
\begin{equation}
 m = \dfrac{1}{c^2} \sqrt{E^2 - P^2 c^2} \approx \qty{137}{\mega\electronvolt}
\end{equation}
La massa ricavata è compatibile con quella del mesone $\pi^0$.
\section{Deutoni}
\textbf{Es. 1}\\
Un fascio di deutoni (D), caratterizzato da una corrente $i_D = \qty{2}{\micro\ampere}$, colpisce un bersaglio di trizio solido con spessore superficiale $\rho d = \qty{0.2}{\milli\gram\per\centi\meter^2}$. Considerando la reazione:
\begin{equation}
 ^3_1\text{H} + ^2_1\text{D} \rightarrow ^4_2\text{He} + n
\end{equation}
Assumendo di avere un ricevoitore con superficie pari a $S=\qty{20}{\centi\meter^2}$ posto ad una distanza $r=\qty{3}{\meter}$ dal bersaglio, calcolare il numero di neutroni al secondo rilevati, sapendo che la sezione d'urto differenziale della reazione è $\dv{\sigma}{\Omega} = \qty{13}{\milli\barn\per\steradian}$ ad un angolo $\theta = 30^\circ$.\\
\textbf{Soluzione:}\\
Calcoliamo il flusso di deutoni che colpisce il bersaglio:
\begin{equation}
 \phi_D = \dfrac{i_D}{q_D S} = \dfrac{i_D}{2e S}
\end{equation}
Il numero di nuclei di trizio nel bersaglio è dato da:
\begin{equation}
 N_{^3\text{H}} = \dfrac{\rho d S N_A}{M_{^3\text{H}}}
\end{equation}
Dove $N_A$ è il numero di Avogadro e $M_{^3\text{H}} \approx \qty{3}{\gram\per\mole}$ è la massa molare del trizio.\\
Poichè l'angolo solido $\Delta \Omega$ è dato da:
\begin{equation}
 \Delta \Omega = \dfrac{S}{r^2}
\end{equation}
e risulta $\Delta \Omega \rightarrow 0 \implies \dv{\sigma}{\Omega} \Delta \Omega \approx \sigma$, possiamo scrivere:
\begin{equation}
  \dot{N}_n = \phi_D N_{^3\text{H}} \dv{\sigma}{\Omega} \Delta \Omega
\end{equation}
Sostituendo i valori numerici, otteniamo:
\begin{equation}
 \dot{N}_n \approx \qty{1.3e4}{\second^{-1}}
\end{equation}
